% Options for packages loaded elsewhere
% Options for packages loaded elsewhere
\PassOptionsToPackage{unicode}{hyperref}
\PassOptionsToPackage{hyphens}{url}
\PassOptionsToPackage{dvipsnames,svgnames,x11names}{xcolor}
%
\documentclass[
  letterpaper,
  DIV=11,
  numbers=noendperiod]{scrartcl}
\usepackage{xcolor}
\usepackage{amsmath,amssymb}
\setcounter{secnumdepth}{-\maxdimen} % remove section numbering
\usepackage{iftex}
\ifPDFTeX
  \usepackage[T1]{fontenc}
  \usepackage[utf8]{inputenc}
  \usepackage{textcomp} % provide euro and other symbols
\else % if luatex or xetex
  \usepackage{unicode-math} % this also loads fontspec
  \defaultfontfeatures{Scale=MatchLowercase}
  \defaultfontfeatures[\rmfamily]{Ligatures=TeX,Scale=1}
\fi
\usepackage{lmodern}
\ifPDFTeX\else
  % xetex/luatex font selection
\fi
% Use upquote if available, for straight quotes in verbatim environments
\IfFileExists{upquote.sty}{\usepackage{upquote}}{}
\IfFileExists{microtype.sty}{% use microtype if available
  \usepackage[]{microtype}
  \UseMicrotypeSet[protrusion]{basicmath} % disable protrusion for tt fonts
}{}
\makeatletter
\@ifundefined{KOMAClassName}{% if non-KOMA class
  \IfFileExists{parskip.sty}{%
    \usepackage{parskip}
  }{% else
    \setlength{\parindent}{0pt}
    \setlength{\parskip}{6pt plus 2pt minus 1pt}}
}{% if KOMA class
  \KOMAoptions{parskip=half}}
\makeatother
% Make \paragraph and \subparagraph free-standing
\makeatletter
\ifx\paragraph\undefined\else
  \let\oldparagraph\paragraph
  \renewcommand{\paragraph}{
    \@ifstar
      \xxxParagraphStar
      \xxxParagraphNoStar
  }
  \newcommand{\xxxParagraphStar}[1]{\oldparagraph*{#1}\mbox{}}
  \newcommand{\xxxParagraphNoStar}[1]{\oldparagraph{#1}\mbox{}}
\fi
\ifx\subparagraph\undefined\else
  \let\oldsubparagraph\subparagraph
  \renewcommand{\subparagraph}{
    \@ifstar
      \xxxSubParagraphStar
      \xxxSubParagraphNoStar
  }
  \newcommand{\xxxSubParagraphStar}[1]{\oldsubparagraph*{#1}\mbox{}}
  \newcommand{\xxxSubParagraphNoStar}[1]{\oldsubparagraph{#1}\mbox{}}
\fi
\makeatother

\usepackage{color}
\usepackage{fancyvrb}
\newcommand{\VerbBar}{|}
\newcommand{\VERB}{\Verb[commandchars=\\\{\}]}
\DefineVerbatimEnvironment{Highlighting}{Verbatim}{commandchars=\\\{\}}
% Add ',fontsize=\small' for more characters per line
\usepackage{framed}
\definecolor{shadecolor}{RGB}{241,243,245}
\newenvironment{Shaded}{\begin{snugshade}}{\end{snugshade}}
\newcommand{\AlertTok}[1]{\textcolor[rgb]{0.68,0.00,0.00}{#1}}
\newcommand{\AnnotationTok}[1]{\textcolor[rgb]{0.37,0.37,0.37}{#1}}
\newcommand{\AttributeTok}[1]{\textcolor[rgb]{0.40,0.45,0.13}{#1}}
\newcommand{\BaseNTok}[1]{\textcolor[rgb]{0.68,0.00,0.00}{#1}}
\newcommand{\BuiltInTok}[1]{\textcolor[rgb]{0.00,0.23,0.31}{#1}}
\newcommand{\CharTok}[1]{\textcolor[rgb]{0.13,0.47,0.30}{#1}}
\newcommand{\CommentTok}[1]{\textcolor[rgb]{0.37,0.37,0.37}{#1}}
\newcommand{\CommentVarTok}[1]{\textcolor[rgb]{0.37,0.37,0.37}{\textit{#1}}}
\newcommand{\ConstantTok}[1]{\textcolor[rgb]{0.56,0.35,0.01}{#1}}
\newcommand{\ControlFlowTok}[1]{\textcolor[rgb]{0.00,0.23,0.31}{\textbf{#1}}}
\newcommand{\DataTypeTok}[1]{\textcolor[rgb]{0.68,0.00,0.00}{#1}}
\newcommand{\DecValTok}[1]{\textcolor[rgb]{0.68,0.00,0.00}{#1}}
\newcommand{\DocumentationTok}[1]{\textcolor[rgb]{0.37,0.37,0.37}{\textit{#1}}}
\newcommand{\ErrorTok}[1]{\textcolor[rgb]{0.68,0.00,0.00}{#1}}
\newcommand{\ExtensionTok}[1]{\textcolor[rgb]{0.00,0.23,0.31}{#1}}
\newcommand{\FloatTok}[1]{\textcolor[rgb]{0.68,0.00,0.00}{#1}}
\newcommand{\FunctionTok}[1]{\textcolor[rgb]{0.28,0.35,0.67}{#1}}
\newcommand{\ImportTok}[1]{\textcolor[rgb]{0.00,0.46,0.62}{#1}}
\newcommand{\InformationTok}[1]{\textcolor[rgb]{0.37,0.37,0.37}{#1}}
\newcommand{\KeywordTok}[1]{\textcolor[rgb]{0.00,0.23,0.31}{\textbf{#1}}}
\newcommand{\NormalTok}[1]{\textcolor[rgb]{0.00,0.23,0.31}{#1}}
\newcommand{\OperatorTok}[1]{\textcolor[rgb]{0.37,0.37,0.37}{#1}}
\newcommand{\OtherTok}[1]{\textcolor[rgb]{0.00,0.23,0.31}{#1}}
\newcommand{\PreprocessorTok}[1]{\textcolor[rgb]{0.68,0.00,0.00}{#1}}
\newcommand{\RegionMarkerTok}[1]{\textcolor[rgb]{0.00,0.23,0.31}{#1}}
\newcommand{\SpecialCharTok}[1]{\textcolor[rgb]{0.37,0.37,0.37}{#1}}
\newcommand{\SpecialStringTok}[1]{\textcolor[rgb]{0.13,0.47,0.30}{#1}}
\newcommand{\StringTok}[1]{\textcolor[rgb]{0.13,0.47,0.30}{#1}}
\newcommand{\VariableTok}[1]{\textcolor[rgb]{0.07,0.07,0.07}{#1}}
\newcommand{\VerbatimStringTok}[1]{\textcolor[rgb]{0.13,0.47,0.30}{#1}}
\newcommand{\WarningTok}[1]{\textcolor[rgb]{0.37,0.37,0.37}{\textit{#1}}}

\usepackage{longtable,booktabs,array}
\usepackage{calc} % for calculating minipage widths
% Correct order of tables after \paragraph or \subparagraph
\usepackage{etoolbox}
\makeatletter
\patchcmd\longtable{\par}{\if@noskipsec\mbox{}\fi\par}{}{}
\makeatother
% Allow footnotes in longtable head/foot
\IfFileExists{footnotehyper.sty}{\usepackage{footnotehyper}}{\usepackage{footnote}}
\makesavenoteenv{longtable}
\usepackage{graphicx}
\makeatletter
\newsavebox\pandoc@box
\newcommand*\pandocbounded[1]{% scales image to fit in text height/width
  \sbox\pandoc@box{#1}%
  \Gscale@div\@tempa{\textheight}{\dimexpr\ht\pandoc@box+\dp\pandoc@box\relax}%
  \Gscale@div\@tempb{\linewidth}{\wd\pandoc@box}%
  \ifdim\@tempb\p@<\@tempa\p@\let\@tempa\@tempb\fi% select the smaller of both
  \ifdim\@tempa\p@<\p@\scalebox{\@tempa}{\usebox\pandoc@box}%
  \else\usebox{\pandoc@box}%
  \fi%
}
% Set default figure placement to htbp
\def\fps@figure{htbp}
\makeatother





\setlength{\emergencystretch}{3em} % prevent overfull lines

\providecommand{\tightlist}{%
  \setlength{\itemsep}{0pt}\setlength{\parskip}{0pt}}



 


\KOMAoption{captions}{tableheading}
\usepackage{amsmath}
\usepackage{amssymb}
\makeatletter
\@ifpackageloaded{tcolorbox}{}{\usepackage[skins,breakable]{tcolorbox}}
\@ifpackageloaded{fontawesome5}{}{\usepackage{fontawesome5}}
\definecolor{quarto-callout-color}{HTML}{909090}
\definecolor{quarto-callout-note-color}{HTML}{0758E5}
\definecolor{quarto-callout-important-color}{HTML}{CC1914}
\definecolor{quarto-callout-warning-color}{HTML}{EB9113}
\definecolor{quarto-callout-tip-color}{HTML}{00A047}
\definecolor{quarto-callout-caution-color}{HTML}{FC5300}
\definecolor{quarto-callout-color-frame}{HTML}{acacac}
\definecolor{quarto-callout-note-color-frame}{HTML}{4582ec}
\definecolor{quarto-callout-important-color-frame}{HTML}{d9534f}
\definecolor{quarto-callout-warning-color-frame}{HTML}{f0ad4e}
\definecolor{quarto-callout-tip-color-frame}{HTML}{02b875}
\definecolor{quarto-callout-caution-color-frame}{HTML}{fd7e14}
\makeatother
\makeatletter
\@ifpackageloaded{caption}{}{\usepackage{caption}}
\AtBeginDocument{%
\ifdefined\contentsname
  \renewcommand*\contentsname{Table of contents}
\else
  \newcommand\contentsname{Table of contents}
\fi
\ifdefined\listfigurename
  \renewcommand*\listfigurename{List of Figures}
\else
  \newcommand\listfigurename{List of Figures}
\fi
\ifdefined\listtablename
  \renewcommand*\listtablename{List of Tables}
\else
  \newcommand\listtablename{List of Tables}
\fi
\ifdefined\figurename
  \renewcommand*\figurename{Figure}
\else
  \newcommand\figurename{Figure}
\fi
\ifdefined\tablename
  \renewcommand*\tablename{Table}
\else
  \newcommand\tablename{Table}
\fi
}
\@ifpackageloaded{float}{}{\usepackage{float}}
\floatstyle{ruled}
\@ifundefined{c@chapter}{\newfloat{codelisting}{h}{lop}}{\newfloat{codelisting}{h}{lop}[chapter]}
\floatname{codelisting}{Listing}
\newcommand*\listoflistings{\listof{codelisting}{List of Listings}}
\makeatother
\makeatletter
\makeatother
\makeatletter
\@ifpackageloaded{caption}{}{\usepackage{caption}}
\@ifpackageloaded{subcaption}{}{\usepackage{subcaption}}
\makeatother
\usepackage{bookmark}
\IfFileExists{xurl.sty}{\usepackage{xurl}}{} % add URL line breaks if available
\urlstyle{same}
\hypersetup{
  pdftitle={Bayesian Neural Networks for Multiclass Classification},
  colorlinks=true,
  linkcolor={blue},
  filecolor={Maroon},
  citecolor={Blue},
  urlcolor={Blue},
  pdfcreator={LaTeX via pandoc}}


\title{Bayesian Neural Networks for Multiclass Classification}
\author{}
\date{}
\begin{document}
\maketitle


\subsection{General Principles}\label{general-principles}

Building upon the binary classification
\href{29.\%20BNN\%20for\%20classification.qmd}{Bayesian Neural Network
(BNN)}, multiclass classification extends the logic to scenarios where
an observation belongs to one of \(K\) mutually exclusive categories
(\(K > 2\)).

Instead of the final layer returning a single output, the final layer in
a multiclass BNN returns a \(K\)-dimensional vector of scores (logits)
for each observation. To transform these continuous scores into valid
probabilities that sum to 1 across all \(K\) classes, we apply the
\textbf{softmax} activation function. Finally, the categorized
predictions are evaluated using a \textbf{Categorical} likelihood.

\subsection{Considerations}\label{considerations}

\begin{tcolorbox}[enhanced jigsaw, colback=white, breakable, left=2mm, arc=.35mm, colbacktitle=quarto-callout-note-color!10!white, colframe=quarto-callout-note-color-frame, leftrule=.75mm, opacityback=0, title=\textcolor{quarto-callout-note-color}{\faInfo}\hspace{0.5em}{Note}, titlerule=0mm, coltitle=black, rightrule=.15mm, toprule=.15mm, bottomtitle=1mm, toptitle=1mm, bottomrule=.15mm, opacitybacktitle=0.6]

\begin{itemize}
\tightlist
\item
  \textbf{Output Layer Dimensions}: While binary classification network
  predictions can be compressed to a single output logit per
  observation, multiclass networks MUST output exactly \(K\) dimensions
  in their final layer, matching the number of target classes.
\item
  \textbf{The Softmax Simplex}: Applying the softmax function across the
  final layer's logits guarantees that the resulting outputs form a
  probability \phantomsection\label{simplex}{{simplex 🛈}}. This is
  biologically similar to independent Poisson rates strictly normalizing
  to fixed categorical ratios.
\item
  \textbf{Likelihood Function}: After calculating the probabilities with
  softmax, we use a \texttt{Categorical} distribution as the final
  likelihood, matching the integer index of the observed category.
\item
  \textbf{Improved Calibration}: Multiclass BNNs greatly reduce
  out-of-distribution overconfidence. Standard deep learning
  cross-entropy models will often assign \textgreater99\% probability to
  an unseen class purely due to the exponential nature of softmax. In a
  BNN, exploring the posterior width of the parameters yields ``flat''
  unconfident probability profiles over \(K\) classes when the input is
  outside the training distribution.
\end{itemize}

\end{tcolorbox}

\subsection{Example}\label{example}

Below is an example code snippet demonstrating a \emph{Bayesian Neural
Network for multiclass classification} using the Bayesian Inference
(\textbf{BI}) package. This example generates a synthetic \(K=3\)
cluster dataset.

\subsection{Python}

\begin{Shaded}
\begin{Highlighting}[]
\ImportTok{from}\NormalTok{ BI }\ImportTok{import}\NormalTok{ bi}
\ImportTok{import}\NormalTok{ jax.numpy }\ImportTok{as}\NormalTok{ jnp}
\ImportTok{import}\NormalTok{ jax}

\CommentTok{\# Setup device{-}{-}{-}{-}{-}{-}{-}{-}{-}{-}{-}{-}{-}{-}{-}{-}{-}{-}{-}{-}{-}{-}{-}{-}{-}{-}{-}{-}{-}{-}{-}{-}{-}{-}{-}{-}{-}{-}{-}{-}{-}{-}{-}{-}{-}{-}{-}{-}}
\NormalTok{m }\OperatorTok{=}\NormalTok{ bi(platform}\OperatorTok{=}\StringTok{\textquotesingle{}cpu\textquotesingle{}}\NormalTok{)}

\CommentTok{\# Generate Synthetic Data {-}{-}{-}{-}{-}{-}{-}{-}{-}{-}{-}{-}{-}{-}{-}{-}{-}{-}{-}{-}{-}{-}{-}{-}{-}{-}{-}{-}{-}{-}{-}{-}{-}{-}{-}{-}}
\CommentTok{\# 3 classes based on a random normal distribution split}
\NormalTok{key }\OperatorTok{=}\NormalTok{ jax.random.PRNGKey(}\DecValTok{42}\NormalTok{)}
\NormalTok{X }\OperatorTok{=}\NormalTok{ jax.random.normal(key, (}\DecValTok{300}\NormalTok{, }\DecValTok{2}\NormalTok{))}
\CommentTok{\# Rule: Q1=Class 0, Q2/Q3=Class 1, Q4=Class 2}
\NormalTok{Y }\OperatorTok{=}\NormalTok{ jnp.where(X[:, }\DecValTok{0}\NormalTok{] }\OperatorTok{\textgreater{}} \DecValTok{0}\NormalTok{, jnp.where(X[:, }\DecValTok{1}\NormalTok{] }\OperatorTok{\textgreater{}} \DecValTok{0}\NormalTok{, }\DecValTok{0}\NormalTok{, }\DecValTok{1}\NormalTok{), }\DecValTok{2}\NormalTok{)}

\NormalTok{m.data\_on\_model }\OperatorTok{=} \BuiltInTok{dict}\NormalTok{(X}\OperatorTok{=}\NormalTok{X, Y}\OperatorTok{=}\NormalTok{Y)}

\CommentTok{\# Define model {-}{-}{-}{-}{-}{-}{-}{-}{-}{-}{-}{-}{-}{-}{-}{-}{-}{-}{-}{-}{-}{-}{-}{-}{-}{-}{-}{-}{-}{-}{-}{-}{-}{-}{-}{-}{-}{-}{-}{-}{-}{-}{-}{-}{-}{-}{-}{-}}
\KeywordTok{def}\NormalTok{ model(X, Y, D\_H1}\OperatorTok{=}\DecValTok{5}\NormalTok{, K}\OperatorTok{=}\DecValTok{3}\NormalTok{):}
\NormalTok{    N, D\_X }\OperatorTok{=}\NormalTok{ X.shape}
    
    \CommentTok{\# First hidden layer: 2 input features {-}\textgreater{} 5 hidden units}
\NormalTok{    w1 }\OperatorTok{=}\NormalTok{ m.bnn.layer\_linear(}
\NormalTok{        X, }
\NormalTok{        dist}\OperatorTok{=}\NormalTok{m.dist.normal(}\DecValTok{0}\NormalTok{, }\DecValTok{1}\NormalTok{, name}\OperatorTok{=}\StringTok{\textquotesingle{}w1\_weight\textquotesingle{}}\NormalTok{, shape}\OperatorTok{=}\NormalTok{(D\_X, D\_H1)),}
\NormalTok{        activation}\OperatorTok{=}\StringTok{\textquotesingle{}tanh\textquotesingle{}}
\NormalTok{    )}
    
    \CommentTok{\# Final output layer: 5 hidden units {-}\textgreater{} K output units}
    \CommentTok{\# Note: No activation is applied automatically inside the layer function here}
\NormalTok{    w2 }\OperatorTok{=}\NormalTok{ m.bnn.layer\_linear(}
\NormalTok{        w1,}
\NormalTok{        dist}\OperatorTok{=}\NormalTok{m.dist.normal(}\DecValTok{0}\NormalTok{, }\DecValTok{1}\NormalTok{, name}\OperatorTok{=}\StringTok{\textquotesingle{}w2\_weight\textquotesingle{}}\NormalTok{, shape}\OperatorTok{=}\NormalTok{(D\_H1, K))}
\NormalTok{    )}
    
    \CommentTok{\# Apply Softmax across the K dimension (axis={-}1) to yield probabilities}
\NormalTok{    p }\OperatorTok{=}\NormalTok{ jax.nn.softmax(w2, axis}\OperatorTok{={-}}\DecValTok{1}\NormalTok{)}
    
    \CommentTok{\# Categorical Likelihood matching indices in Y}
\NormalTok{    m.dist.categorical(probs}\OperatorTok{=}\NormalTok{p, obs}\OperatorTok{=}\NormalTok{Y)}

\CommentTok{\# Run mcmc {-}{-}{-}{-}{-}{-}{-}{-}{-}{-}{-}{-}{-}{-}{-}{-}{-}{-}{-}{-}{-}{-}{-}{-}{-}{-}{-}{-}{-}{-}{-}{-}{-}{-}{-}{-}{-}{-}{-}{-}{-}{-}{-}{-}{-}{-}{-}{-}}
\NormalTok{m.fit(model, num\_samples}\OperatorTok{=}\DecValTok{500}\NormalTok{, progress\_bar}\OperatorTok{=}\VariableTok{False}\NormalTok{) }\CommentTok{\# Approximate posterior distributions}

\CommentTok{\# Predictions from the model {-}{-}{-}{-}{-}{-}{-}{-}{-}{-}{-}{-}{-}{-}{-}{-}{-}{-}{-}{-}{-}{-}{-}{-}{-}{-}{-}{-}{-}{-}{-}{-}{-}{-}{-}{-}{-}{-}{-}{-}{-}{-}{-}{-}{-}{-}{-}{-}}
\ImportTok{import}\NormalTok{ matplotlib.pyplot }\ImportTok{as}\NormalTok{ plt}

\CommentTok{\# Create a grid to evaluate the model}
\NormalTok{n\_grid }\OperatorTok{=} \DecValTok{50}
\NormalTok{x0 }\OperatorTok{=}\NormalTok{ jnp.linspace(X[:, }\DecValTok{0}\NormalTok{].}\BuiltInTok{min}\NormalTok{() }\OperatorTok{{-}} \FloatTok{0.5}\NormalTok{, X[:, }\DecValTok{0}\NormalTok{].}\BuiltInTok{max}\NormalTok{() }\OperatorTok{+} \FloatTok{0.5}\NormalTok{, n\_grid)}
\NormalTok{x1 }\OperatorTok{=}\NormalTok{ jnp.linspace(X[:, }\DecValTok{1}\NormalTok{].}\BuiltInTok{min}\NormalTok{() }\OperatorTok{{-}} \FloatTok{0.5}\NormalTok{, X[:, }\DecValTok{1}\NormalTok{].}\BuiltInTok{max}\NormalTok{() }\OperatorTok{+} \FloatTok{0.5}\NormalTok{, n\_grid)}
\NormalTok{xx0, xx1 }\OperatorTok{=}\NormalTok{ jnp.meshgrid(x0, x1)}
\NormalTok{X\_grid }\OperatorTok{=}\NormalTok{ jnp.c\_[xx0.ravel(), xx1.ravel()]}

\CommentTok{\# Swap data on model temporarily to predict on the grid}
\NormalTok{m.data\_on\_model }\OperatorTok{=} \BuiltInTok{dict}\NormalTok{(X}\OperatorTok{=}\NormalTok{X\_grid, Y}\OperatorTok{=}\NormalTok{jnp.zeros(X\_grid.shape[}\DecValTok{0}\NormalTok{], dtype}\OperatorTok{=}\NormalTok{jnp.int32))}
\NormalTok{pred }\OperatorTok{=}\NormalTok{ m.sample(samples}\OperatorTok{=}\DecValTok{500}\NormalTok{)[}\StringTok{\textquotesingle{}Y\textquotesingle{}}\NormalTok{]}
\NormalTok{p\_mean }\OperatorTok{=}\NormalTok{ jnp.mean(pred, axis}\OperatorTok{=}\DecValTok{0}\NormalTok{)}

\CommentTok{\# Plotting the posterior predictive mean (categorical blending)}
\NormalTok{fig, ax }\OperatorTok{=}\NormalTok{ plt.subplots(figsize}\OperatorTok{=}\NormalTok{(}\DecValTok{8}\NormalTok{, }\DecValTok{6}\NormalTok{), constrained\_layout}\OperatorTok{=}\VariableTok{True}\NormalTok{)}
\NormalTok{contour }\OperatorTok{=}\NormalTok{ ax.contourf(xx0, xx1, p\_mean.reshape(n\_grid, n\_grid), cmap}\OperatorTok{=}\StringTok{"viridis"}\NormalTok{, alpha}\OperatorTok{=}\FloatTok{0.6}\NormalTok{)}
\NormalTok{scatter }\OperatorTok{=}\NormalTok{ ax.scatter(X[:, }\DecValTok{0}\NormalTok{], X[:, }\DecValTok{1}\NormalTok{], c}\OperatorTok{=}\NormalTok{Y, cmap}\OperatorTok{=}\StringTok{"viridis"}\NormalTok{, edgecolors}\OperatorTok{=}\StringTok{\textquotesingle{}k\textquotesingle{}}\NormalTok{)}
\NormalTok{ax.}\BuiltInTok{set}\NormalTok{(title}\OperatorTok{=}\StringTok{"Posterior Predictive Mean"}\NormalTok{, xlabel}\OperatorTok{=}\StringTok{"Feature 1"}\NormalTok{, ylabel}\OperatorTok{=}\StringTok{"Feature 2"}\NormalTok{)}
\NormalTok{fig.colorbar(contour, ax}\OperatorTok{=}\NormalTok{ax)}
\end{Highlighting}
\end{Shaded}

\subsection{Julia}

\begin{Shaded}
\begin{Highlighting}[]
\ImportTok{using} \BuiltInTok{BayesianInference}
\ImportTok{using} \BuiltInTok{PythonCall}

\CommentTok{\# Setup device{-}{-}{-}{-}{-}{-}{-}{-}{-}{-}{-}{-}{-}{-}{-}{-}{-}{-}{-}{-}{-}{-}{-}{-}{-}{-}{-}{-}{-}{-}{-}{-}{-}{-}{-}{-}{-}{-}{-}{-}{-}{-}{-}{-}{-}{-}{-}{-}}
\NormalTok{m }\OperatorTok{=} \FunctionTok{importBI}\NormalTok{(platform}\OperatorTok{=}\StringTok{"cpu"}\NormalTok{)}

\CommentTok{\# Generate Synthetic Data {-}{-}{-}{-}{-}{-}{-}{-}{-}{-}{-}{-}{-}{-}{-}{-}{-}{-}{-}{-}{-}{-}{-}{-}{-}{-}{-}{-}{-}{-}{-}{-}{-}{-}{-}{-}}
\NormalTok{np }\OperatorTok{=} \FunctionTok{pyimport}\NormalTok{(}\StringTok{"numpy"}\NormalTok{)}
\NormalTok{jax\_random }\OperatorTok{=} \FunctionTok{pyimport}\NormalTok{(}\StringTok{"jax.random"}\NormalTok{)}
\NormalTok{jnp }\OperatorTok{=} \FunctionTok{pyimport}\NormalTok{(}\StringTok{"jax.numpy"}\NormalTok{)}

\NormalTok{key }\OperatorTok{=}\NormalTok{ jax\_random.}\FunctionTok{PRNGKey}\NormalTok{(}\FloatTok{42}\NormalTok{)}
\NormalTok{X }\OperatorTok{=}\NormalTok{ jax\_random.}\FunctionTok{normal}\NormalTok{(key, (}\FloatTok{300}\NormalTok{, }\FloatTok{2}\NormalTok{))}

\CommentTok{\# Simple rule to partition into K=3 classes}
\NormalTok{Y }\OperatorTok{=}\NormalTok{ jnp.}\FunctionTok{where}\NormalTok{(X[}\OperatorTok{:}\NormalTok{, }\FloatTok{0}\NormalTok{] }\OperatorTok{\textgreater{}} \FloatTok{0}\NormalTok{, jnp.}\FunctionTok{where}\NormalTok{(X[}\OperatorTok{:}\NormalTok{, }\FloatTok{1}\NormalTok{] }\OperatorTok{\textgreater{}} \FloatTok{0}\NormalTok{, }\FloatTok{0}\NormalTok{, }\FloatTok{1}\NormalTok{), }\FloatTok{2}\NormalTok{)}

\NormalTok{m.data\_on\_model[}\StringTok{"X"}\NormalTok{] }\OperatorTok{=}\NormalTok{ X}
\NormalTok{m.data\_on\_model[}\StringTok{"Y"}\NormalTok{] }\OperatorTok{=}\NormalTok{ Y}

\CommentTok{\# Define model {-}{-}{-}{-}{-}{-}{-}{-}{-}{-}{-}{-}{-}{-}{-}{-}{-}{-}{-}{-}{-}{-}{-}{-}{-}{-}{-}{-}{-}{-}{-}{-}{-}{-}{-}{-}{-}{-}{-}{-}{-}{-}{-}{-}{-}{-}{-}{-}}
\PreprocessorTok{@BI} \KeywordTok{function} \FunctionTok{model}\NormalTok{(X, Y)}
\NormalTok{    N, D\_X }\OperatorTok{=} \FunctionTok{size}\NormalTok{(X)}
\NormalTok{    D\_H1 }\OperatorTok{=} \FloatTok{5}
\NormalTok{    K }\OperatorTok{=} \FloatTok{3}
    
    \CommentTok{\# First hidden layer}
\NormalTok{    w1 }\OperatorTok{=}\NormalTok{ m.bnn.}\FunctionTok{layer\_linear}\NormalTok{(}
\NormalTok{        X, }
\NormalTok{        dist}\OperatorTok{=}\NormalTok{m.dist.}\FunctionTok{normal}\NormalTok{(}\FloatTok{0}\NormalTok{, }\FloatTok{1}\NormalTok{, name}\OperatorTok{=}\StringTok{"w1\_weight"}\NormalTok{, shape}\OperatorTok{=}\NormalTok{(D\_X, D\_H1)),}
\NormalTok{        activation}\OperatorTok{=}\StringTok{"tanh"}
\NormalTok{    )}
    
    \CommentTok{\# Final output layer }
\NormalTok{    w2 }\OperatorTok{=}\NormalTok{ m.bnn.}\FunctionTok{layer\_linear}\NormalTok{(}
\NormalTok{        w1,}
\NormalTok{        dist}\OperatorTok{=}\NormalTok{m.dist.}\FunctionTok{normal}\NormalTok{(}\FloatTok{0}\NormalTok{, }\FloatTok{1}\NormalTok{, name}\OperatorTok{=}\StringTok{"w2\_weight"}\NormalTok{, shape}\OperatorTok{=}\NormalTok{(D\_H1, K))}
\NormalTok{    )}
    
    \CommentTok{\# Softmax conversion to probability simplex}
\NormalTok{    p }\OperatorTok{=}\NormalTok{ jax.nn.}\FunctionTok{softmax}\NormalTok{(w2, axis}\OperatorTok{={-}}\FloatTok{1}\NormalTok{)}
    
    \CommentTok{\# Categorical Likelihood}
\NormalTok{    m.dist.}\FunctionTok{categorical}\NormalTok{(probs}\OperatorTok{=}\NormalTok{p, obs}\OperatorTok{=}\NormalTok{Y)}
\KeywordTok{end}

\CommentTok{\# Run mcmc {-}{-}{-}{-}{-}{-}{-}{-}{-}{-}{-}{-}{-}{-}{-}{-}{-}{-}{-}{-}{-}{-}{-}{-}{-}{-}{-}{-}{-}{-}{-}{-}{-}{-}{-}{-}{-}{-}{-}{-}{-}{-}{-}{-}{-}{-}{-}{-}}
\NormalTok{m.}\FunctionTok{fit}\NormalTok{(model, num\_samples}\OperatorTok{=}\FloatTok{500}\NormalTok{, progress\_bar}\OperatorTok{=}\ConstantTok{false}\NormalTok{)}
\end{Highlighting}
\end{Shaded}

\subsection{Mathematical Details}\label{mathematical-details}

\subsubsection{\texorpdfstring{\emph{Bayesian
Formulation}}{Bayesian Formulation}}\label{bayesian-formulation}

For a multiclass classification task spanning \(N\) observations and
\(K\) mutually exclusive classes, we model the probability vector
\(\theta_i\) that the response \(Y_i \in \{0, 1, ..., K-1\}\) falls into
each respective class.

\begin{enumerate}
\def\labelenumi{\arabic{enumi}.}
\item
  \textbf{Hidden Layer}: \[
  \Omega_{i,1} = \tanh(X_i \Theta_1)
  \] Where \(X_i\) is a \(2\)-vector of features and \(\Theta_1\) is a
  \(2 \times 5\) weight matrix parameterized by a Standard Normal prior:
  \(\Theta_1 \sim \text{Normal}(0, 1)\).
\item
  \textbf{Output Layer (Logits)}: \[
  \phi_i = \Omega_{i,1} \Theta_2
  \] Where \(\Theta_2\) is a \(5 \times 3\) weight matrix mapping the
  hidden features to the logits for the \(K=3\) classes. This is
  parameterized by: \(\Theta_2 \sim \text{Normal}(0, 1)\).
\item
  \textbf{Softmax Transformation}: \[
  \theta_i = \text{Softmax}(\phi_i)
  \] The softmax function ensures that for every observation \(i\), the
  sum of probabilities across the 3 classes equals exactly 1.
\item
  \textbf{Likelihood}: \[
  Y_i \sim \text{Categorical}(\theta_i)
  \] The observed class index \(Y_i\) is drawn from a Categorical
  distribution parameterized by the probability vector \(\theta_i\).
\end{enumerate}

\subsection{Notes}\label{notes}

\begin{tcolorbox}[enhanced jigsaw, colback=white, breakable, left=2mm, arc=.35mm, colbacktitle=quarto-callout-note-color!10!white, colframe=quarto-callout-note-color-frame, leftrule=.75mm, opacityback=0, title=\textcolor{quarto-callout-note-color}{\faInfo}\hspace{0.5em}{Note}, titlerule=0mm, coltitle=black, rightrule=.15mm, toprule=.15mm, bottomtitle=1mm, toptitle=1mm, bottomrule=.15mm, opacitybacktitle=0.6]

\begin{itemize}
\tightlist
\item
  For large outputs where \(K > 100\), computing the exact softmax
  normalization scalar (the denominator term combining all exponentiated
  logits) can become computationally expensive over thousands of MCMC
  posterior evaluations.
\item
  Neural networks configured with a standard Cross-Entropy loss mapping
  to one-hot vectors conceptually perform exactly this sequence: dot
  product of final weights \(\rightarrow\) Softmax \(\rightarrow\)
  Categorical Likelihood.
\end{itemize}

\end{tcolorbox}

\subsection{Reference(s)}\label{references}

\begin{enumerate}
\def\labelenumi{\arabic{enumi}.}
\tightlist
\item
  \href{https://juanitorduz.github.io/intro_svi/}{PyData Berlin 2025:
  Introduction to Stochastic Variational Inference with NumPyro}
\end{enumerate}




\end{document}
